% Options for packages loaded elsewhere
\PassOptionsToPackage{unicode}{hyperref}
\PassOptionsToPackage{hyphens}{url}
\documentclass[
]{article}
\usepackage{xcolor}
\usepackage[margin=1in]{geometry}
\usepackage{amsmath,amssymb}
\setcounter{secnumdepth}{-\maxdimen} % remove section numbering
\usepackage{iftex}
\ifPDFTeX
  \usepackage[T1]{fontenc}
  \usepackage[utf8]{inputenc}
  \usepackage{textcomp} % provide euro and other symbols
\else % if luatex or xetex
  \usepackage{unicode-math} % this also loads fontspec
  \defaultfontfeatures{Scale=MatchLowercase}
  \defaultfontfeatures[\rmfamily]{Ligatures=TeX,Scale=1}
\fi
\usepackage{lmodern}
\ifPDFTeX\else
  % xetex/luatex font selection
\fi
% Use upquote if available, for straight quotes in verbatim environments
\IfFileExists{upquote.sty}{\usepackage{upquote}}{}
\IfFileExists{microtype.sty}{% use microtype if available
  \usepackage[]{microtype}
  \UseMicrotypeSet[protrusion]{basicmath} % disable protrusion for tt fonts
}{}
\makeatletter
\@ifundefined{KOMAClassName}{% if non-KOMA class
  \IfFileExists{parskip.sty}{%
    \usepackage{parskip}
  }{% else
    \setlength{\parindent}{0pt}
    \setlength{\parskip}{6pt plus 2pt minus 1pt}}
}{% if KOMA class
  \KOMAoptions{parskip=half}}
\makeatother
\usepackage{color}
\usepackage{fancyvrb}
\newcommand{\VerbBar}{|}
\newcommand{\VERB}{\Verb[commandchars=\\\{\}]}
\DefineVerbatimEnvironment{Highlighting}{Verbatim}{commandchars=\\\{\}}
% Add ',fontsize=\small' for more characters per line
\usepackage{framed}
\definecolor{shadecolor}{RGB}{248,248,248}
\newenvironment{Shaded}{\begin{snugshade}}{\end{snugshade}}
\newcommand{\AlertTok}[1]{\textcolor[rgb]{0.94,0.16,0.16}{#1}}
\newcommand{\AnnotationTok}[1]{\textcolor[rgb]{0.56,0.35,0.01}{\textbf{\textit{#1}}}}
\newcommand{\AttributeTok}[1]{\textcolor[rgb]{0.13,0.29,0.53}{#1}}
\newcommand{\BaseNTok}[1]{\textcolor[rgb]{0.00,0.00,0.81}{#1}}
\newcommand{\BuiltInTok}[1]{#1}
\newcommand{\CharTok}[1]{\textcolor[rgb]{0.31,0.60,0.02}{#1}}
\newcommand{\CommentTok}[1]{\textcolor[rgb]{0.56,0.35,0.01}{\textit{#1}}}
\newcommand{\CommentVarTok}[1]{\textcolor[rgb]{0.56,0.35,0.01}{\textbf{\textit{#1}}}}
\newcommand{\ConstantTok}[1]{\textcolor[rgb]{0.56,0.35,0.01}{#1}}
\newcommand{\ControlFlowTok}[1]{\textcolor[rgb]{0.13,0.29,0.53}{\textbf{#1}}}
\newcommand{\DataTypeTok}[1]{\textcolor[rgb]{0.13,0.29,0.53}{#1}}
\newcommand{\DecValTok}[1]{\textcolor[rgb]{0.00,0.00,0.81}{#1}}
\newcommand{\DocumentationTok}[1]{\textcolor[rgb]{0.56,0.35,0.01}{\textbf{\textit{#1}}}}
\newcommand{\ErrorTok}[1]{\textcolor[rgb]{0.64,0.00,0.00}{\textbf{#1}}}
\newcommand{\ExtensionTok}[1]{#1}
\newcommand{\FloatTok}[1]{\textcolor[rgb]{0.00,0.00,0.81}{#1}}
\newcommand{\FunctionTok}[1]{\textcolor[rgb]{0.13,0.29,0.53}{\textbf{#1}}}
\newcommand{\ImportTok}[1]{#1}
\newcommand{\InformationTok}[1]{\textcolor[rgb]{0.56,0.35,0.01}{\textbf{\textit{#1}}}}
\newcommand{\KeywordTok}[1]{\textcolor[rgb]{0.13,0.29,0.53}{\textbf{#1}}}
\newcommand{\NormalTok}[1]{#1}
\newcommand{\OperatorTok}[1]{\textcolor[rgb]{0.81,0.36,0.00}{\textbf{#1}}}
\newcommand{\OtherTok}[1]{\textcolor[rgb]{0.56,0.35,0.01}{#1}}
\newcommand{\PreprocessorTok}[1]{\textcolor[rgb]{0.56,0.35,0.01}{\textit{#1}}}
\newcommand{\RegionMarkerTok}[1]{#1}
\newcommand{\SpecialCharTok}[1]{\textcolor[rgb]{0.81,0.36,0.00}{\textbf{#1}}}
\newcommand{\SpecialStringTok}[1]{\textcolor[rgb]{0.31,0.60,0.02}{#1}}
\newcommand{\StringTok}[1]{\textcolor[rgb]{0.31,0.60,0.02}{#1}}
\newcommand{\VariableTok}[1]{\textcolor[rgb]{0.00,0.00,0.00}{#1}}
\newcommand{\VerbatimStringTok}[1]{\textcolor[rgb]{0.31,0.60,0.02}{#1}}
\newcommand{\WarningTok}[1]{\textcolor[rgb]{0.56,0.35,0.01}{\textbf{\textit{#1}}}}
\usepackage{graphicx}
\makeatletter
\newsavebox\pandoc@box
\newcommand*\pandocbounded[1]{% scales image to fit in text height/width
  \sbox\pandoc@box{#1}%
  \Gscale@div\@tempa{\textheight}{\dimexpr\ht\pandoc@box+\dp\pandoc@box\relax}%
  \Gscale@div\@tempb{\linewidth}{\wd\pandoc@box}%
  \ifdim\@tempb\p@<\@tempa\p@\let\@tempa\@tempb\fi% select the smaller of both
  \ifdim\@tempa\p@<\p@\scalebox{\@tempa}{\usebox\pandoc@box}%
  \else\usebox{\pandoc@box}%
  \fi%
}
% Set default figure placement to htbp
\def\fps@figure{htbp}
\makeatother
\setlength{\emergencystretch}{3em} % prevent overfull lines
\providecommand{\tightlist}{%
  \setlength{\itemsep}{0pt}\setlength{\parskip}{0pt}}
\usepackage{bookmark}
\IfFileExists{xurl.sty}{\usepackage{xurl}}{} % add URL line breaks if available
\urlstyle{same}
\hypersetup{
  pdftitle={Midterm 2 W26},
  hidelinks,
  pdfcreator={LaTeX via pandoc}}

\title{Midterm 2 W26}
\author{}
\date{\vspace{-2.5em}2026-03-03}

\begin{document}
\maketitle

\subsection{Instructions}\label{instructions}

This is an in-person exam. Your exam must be complete and uploaded to
Gradescope by 2:00p.

Answer the following questions and complete the exercises in RMarkdown.
Please embed all of your code and push your final work to your
repository. Your code must be organized, clean, and run free from
errors. Remember, you must remove the \texttt{\#} for any included code
chunks to run. Be sure to add your name to the author header above.

Your code must knit in order to be considered. If you are stuck and
cannot answer a question, then comment out your code and knit the
document. You may use your notes, labs, and homework to help you
complete this exam. Do not use any other resources- including AI
assistance or other students' work.

Don't forget to answer any questions that are asked in the prompt! Each
question must be coded; it cannot be answered by a sort in a spreadsheet
or a written response only.

For all plots you create, a title and clearly labeled axes must be
provided.

Be sure to push your completed midterm to your repository and upload the
document to Gradescope. \emph{Incorrectly uploading the exam to
Gradescope will result in a 5 point deduction.} This exam is worth 50
points.

\subsection{Load the libraries}\label{load-the-libraries}

\begin{Shaded}
\begin{Highlighting}[]
\FunctionTok{library}\NormalTok{(tidyverse)}
\FunctionTok{library}\NormalTok{(janitor)}
\FunctionTok{library}\NormalTok{(naniar)}
\FunctionTok{library}\NormalTok{(ggmap)}
\FunctionTok{library}\NormalTok{(lubridate)}
\FunctionTok{library}\NormalTok{(shiny)}
\end{Highlighting}
\end{Shaded}

\subsection{API Key}\label{api-key}

Since you will build a map as part of the exam, take a moment to
register your API key. You may need to look up your
\href{https://client.stadiamaps.com/dashboard/overview}{stadia maps} API
key online (please do that now).

\begin{Shaded}
\begin{Highlighting}[]
\FunctionTok{register\_stadiamaps}\NormalTok{(}\StringTok{"7e1ebae2{-}6056{-}40c4{-}a2af{-}d1b197485428"}\NormalTok{, }\AttributeTok{write =} \ConstantTok{FALSE}\NormalTok{)}
\end{Highlighting}
\end{Shaded}

\subsection{About the data}\label{about-the-data}

This study is focused on the fasting behavior of polar bears in the
Southern Beaufort (SB) and Chukchi seas (CS). The data were collected
from 1983-2017 and include measurements of blood urea nitrogen (BUN) and
creatinine, which are used to calculate a ratio that indicates whether a
bear is fasting or not. The original paper is included in the data
folder. \emph{It may help to quickly read the abstract before starting
the exam.}

The data are from: Rode, K.D., Wilson, R.R., Douglas, D.C., Muhlenbruch,
V., Atwood, T.C., Regehr, E.V., Richardson, E.S., Pilfold, N.W.,
Derocher, A.E., Durner, G.M., Stirling, I., Amstrup, S.C., St.~Martin,
M., Pagano, A.M., and others, 2017, Spring fasting behavior in a marine
apex predator provides an index of ecosystem productivity: Global Change
Biology, v. 24, no. 1, p.~410--423,
\url{https://doi.org/10.1111/gcb.13933}.

\subsection{Data resources}\label{data-resources}

\begin{itemize}
\tightlist
\item
  \url{https://data.usgs.gov/datacatalog/data/USGS:ASC141}
\item
  \url{https://www.sciencebase.gov/catalog/item/61a1a9e7d34eb622f69743d2}
\end{itemize}

\subsection{The variables include:}\label{the-variables-include}

\begin{itemize}
\tightlist
\item
  \texttt{bear\_id}: A unique number given to each bear.\\
\item
  \texttt{capture\_date}: The data when the bear was tranquilized and
  assessed.\\
\item
  \texttt{popn}: The population to which the bear belongs.\\
\item
  \texttt{sex}: Sex of each bear.\\
\item
  \texttt{age\_class}: The age class of each bear (A=adult,
  S=subadult).\\
\item
  \texttt{mating}: Whether or not the bear is pursuing a mate.\\
\item
  \texttt{denning}: Whether or not the bear denned last year.\\
\item
  \texttt{lat}: Latitude where the bear was located.\\
\item
  \texttt{long}: Longitude where the bear was located.\\
\item
  \texttt{urea\_nitrogen}: Concentration of urea nitrogen.\\
\item
  \texttt{creatinine}: Concentration of creatinine.
\end{itemize}

\subsection{Load the data}\label{load-the-data}

Run the code chunk below to load the data, clean the names, and do some
tidying.

\begin{Shaded}
\begin{Highlighting}[]
\NormalTok{polar\_bears }\OtherTok{\textless{}{-}} 
  \FunctionTok{read\_csv}\NormalTok{(}\StringTok{"data/polarBear\_serumUreaCreatinine\_beaufortChukchi\_rode.csv"}\NormalTok{) }\SpecialCharTok{\%\textgreater{}\%} 
  \FunctionTok{clean\_names}\NormalTok{() }\SpecialCharTok{\%\textgreater{}\%} 
  \FunctionTok{select}\NormalTok{(}\SpecialCharTok{{-}}\NormalTok{den\_method, }\SpecialCharTok{{-}}\NormalTok{no\_cubs, }\SpecialCharTok{{-}}\NormalTok{cub\_age) }\SpecialCharTok{\%\textgreater{}\%}
  \FunctionTok{filter}\NormalTok{(age\_class}\SpecialCharTok{==}\FunctionTok{c}\NormalTok{(}\StringTok{"A"}\NormalTok{, }\StringTok{"S"}\NormalTok{)) }\SpecialCharTok{\%\textgreater{}\%} 
  \FunctionTok{mutate}\NormalTok{(}\AttributeTok{sex=}\FunctionTok{toupper}\NormalTok{(sex), }
         \AttributeTok{age\_class=}\FunctionTok{toupper}\NormalTok{(age\_class))}
\end{Highlighting}
\end{Shaded}

\begin{enumerate}
\def\labelenumi{\arabic{enumi}.}
\tightlist
\item
  (2 points) List the variable names and use a summary function to get
  an idea of the data structure.
\end{enumerate}

\begin{Shaded}
\begin{Highlighting}[]
\FunctionTok{summary}\NormalTok{(polar\_bears)}
\end{Highlighting}
\end{Shaded}

\begin{verbatim}
##     bear_id       capture_date            popn               sex           
##  Min.   :  697   Min.   :1983-04-23   Length:519         Length:519        
##  1st Qu.:20166   1st Qu.:2000-04-10   Class :character   Class :character  
##  Median :20705   Median :2008-04-12   Mode  :character   Mode  :character  
##  Mean   :18040   Mean   :2004-07-02                                        
##  3rd Qu.:21220   3rd Qu.:2011-03-31                                        
##  Max.   :99990   Max.   :2016-04-22                                        
##   age_class            mating            denning               lat       
##  Length:519         Length:519         Length:519         Min.   :62.46  
##  Class :character   Class :character   Class :character   1st Qu.:68.11  
##  Mode  :character   Mode  :character   Mode  :character   Median :70.53  
##                                                           Mean   :69.86  
##                                                           3rd Qu.:71.00  
##                                                           Max.   :72.70  
##       long        urea_nitrogen     creatinine   
##  Min.   :-172.2   Min.   : 1.00   Min.   :0.350  
##  1st Qu.:-165.2   1st Qu.: 6.00   1st Qu.:0.850  
##  Median :-150.7   Median :11.00   Median :1.050  
##  Mean   :-153.7   Mean   :15.13   Mean   :1.079  
##  3rd Qu.:-146.5   3rd Qu.:21.25   3rd Qu.:1.285  
##  Max.   :-134.8   Max.   :69.00   Max.   :2.350
\end{verbatim}

\begin{Shaded}
\begin{Highlighting}[]
\FunctionTok{names}\NormalTok{(polar\_bears)}
\end{Highlighting}
\end{Shaded}

\begin{verbatim}
##  [1] "bear_id"       "capture_date"  "popn"          "sex"          
##  [5] "age_class"     "mating"        "denning"       "lat"          
##  [9] "long"          "urea_nitrogen" "creatinine"
\end{verbatim}

\begin{enumerate}
\def\labelenumi{\arabic{enumi}.}
\setcounter{enumi}{1}
\tightlist
\item
  (3 points) How are NA's represented and where are they distributed in
  the variables? Use the \texttt{naniar} package to visualize the
  distribution of NA's in the data and replace any alternative
  representations of NA with actual NA values.
\end{enumerate}

\begin{Shaded}
\begin{Highlighting}[]
\NormalTok{polar\_bears}
\end{Highlighting}
\end{Shaded}

\begin{verbatim}
## # A tibble: 519 x 11
##    bear_id capture_date popn  sex   age_class mating denning   lat  long
##      <dbl> <date>       <chr> <chr> <chr>     <chr>  <chr>   <dbl> <dbl>
##  1    6398 1985-04-28   SB    F     A         no     yes      70.4 -148.
##  2    6384 1991-04-13   SB    F     A         no     yes      70.0 -150.
##  3    6520 1992-04-18   SB    F     A         no     yes      70.9 -145.
##  4    6559 2000-04-12   SB    F     A         no     yes      70.0 -142.
##  5   20629 2003-04-16   SB    F     A         no     yes      70.5 -145.
##  6   20876 2007-04-01   SB    F     A         no     yes      71.1 -153.
##  7   20354 2009-04-16   SB    F     A         no     yes      70.7 -145.
##  8   21024 2009-04-25   SB    F     A         no     yes      70.5 -147.
##  9   20764 2009-05-01   SB    F     A         no     yes      70.8 -149.
## 10   21045 2009-05-19   SB    F     A         no     yes      70.7 -148.
## # i 509 more rows
## # i 2 more variables: urea_nitrogen <dbl>, creatinine <dbl>
\end{verbatim}

\begin{Shaded}
\begin{Highlighting}[]
\CommentTok{\#NAs represented as unknown , in mating}
\end{Highlighting}
\end{Shaded}

\begin{Shaded}
\begin{Highlighting}[]
\NormalTok{polar\_bears\_no\_nas }\OtherTok{\textless{}{-}} 
  \FunctionTok{read\_csv}\NormalTok{(}\StringTok{"data/polarBear\_serumUreaCreatinine\_beaufortChukchi\_rode.csv"}\NormalTok{, }\AttributeTok{na=}\StringTok{"{-}999"}\NormalTok{) }\SpecialCharTok{\%\textgreater{}\%} 
  \FunctionTok{clean\_names}\NormalTok{()}
\end{Highlighting}
\end{Shaded}

\begin{verbatim}
## Rows: 1048 Columns: 14
## -- Column specification --------------------------------------------------------
## Delimiter: ","
## chr  (8): Popn, Sex, AgeClass, Mating, Denning, den_method, NoCubs, CubAge
## dbl  (5): BearID, Lat, Long, Urea Nitrogen, Creatinine
## date (1): Capture Date
## 
## i Use `spec()` to retrieve the full column specification for this data.
## i Specify the column types or set `show_col_types = FALSE` to quiet this message.
\end{verbatim}

\begin{Shaded}
\begin{Highlighting}[]
\NormalTok{polar\_bears }\SpecialCharTok{\%\textgreater{}\%} 
  \FunctionTok{summarize}\NormalTok{(}\FunctionTok{across}\NormalTok{(}\FunctionTok{everything}\NormalTok{(), }\SpecialCharTok{\textasciitilde{}} \FunctionTok{sum}\NormalTok{(}\FunctionTok{is.na}\NormalTok{(.))))}
\end{Highlighting}
\end{Shaded}

\begin{verbatim}
## # A tibble: 1 x 11
##   bear_id capture_date  popn   sex age_class mating denning   lat  long
##     <int>        <int> <int> <int>     <int>  <int>   <int> <int> <int>
## 1       0            0     0     0         0      0       0     0     0
## # i 2 more variables: urea_nitrogen <int>, creatinine <int>
\end{verbatim}

\begin{Shaded}
\begin{Highlighting}[]
\FunctionTok{miss\_var\_summary}\NormalTok{(polar\_bears\_no\_nas)}
\end{Highlighting}
\end{Shaded}

\begin{verbatim}
## # A tibble: 14 x 3
##    variable      n_miss pct_miss
##    <chr>          <int>    <num>
##  1 bear_id            0        0
##  2 capture_date       0        0
##  3 popn               0        0
##  4 sex                0        0
##  5 age_class          0        0
##  6 mating             0        0
##  7 denning            0        0
##  8 den_method         0        0
##  9 lat                0        0
## 10 long               0        0
## 11 no_cubs            0        0
## 12 cub_age            0        0
## 13 urea_nitrogen      0        0
## 14 creatinine         0        0
\end{verbatim}

\begin{Shaded}
\begin{Highlighting}[]
\NormalTok{polar\_bears\_no\_nas }\SpecialCharTok{\%\textgreater{}\%} 
  \FunctionTok{gg\_miss\_var}\NormalTok{()}
\end{Highlighting}
\end{Shaded}

\pandocbounded{\includegraphics[keepaspectratio]{midterm_2_files/figure-latex/unnamed-chunk-9-1.pdf}}

\begin{enumerate}
\def\labelenumi{\arabic{enumi}.}
\setcounter{enumi}{2}
\tightlist
\item
  (6 points) Is sampling even for both sexes? Make a count the shows the
  number of males and females sampled by population and make a plot to
  visualize the counts.
\end{enumerate}

\begin{Shaded}
\begin{Highlighting}[]
\FunctionTok{names}\NormalTok{(polar\_bears)}
\end{Highlighting}
\end{Shaded}

\begin{verbatim}
##  [1] "bear_id"       "capture_date"  "popn"          "sex"          
##  [5] "age_class"     "mating"        "denning"       "lat"          
##  [9] "long"          "urea_nitrogen" "creatinine"
\end{verbatim}

\begin{Shaded}
\begin{Highlighting}[]
\NormalTok{polar\_bears }\SpecialCharTok{\%\textgreater{}\%} 
  \FunctionTok{count}\NormalTok{(sex, }\AttributeTok{sort=}\NormalTok{T)}
\end{Highlighting}
\end{Shaded}

\begin{verbatim}
## # A tibble: 2 x 2
##   sex       n
##   <chr> <int>
## 1 F       292
## 2 M       227
\end{verbatim}

\begin{Shaded}
\begin{Highlighting}[]
\NormalTok{polar\_bears }\SpecialCharTok{\%\textgreater{}\%} 
  \FunctionTok{ggplot}\NormalTok{(}\FunctionTok{aes}\NormalTok{(}\AttributeTok{x=}\NormalTok{sex))}\SpecialCharTok{+}
  \FunctionTok{geom\_bar}\NormalTok{()}\SpecialCharTok{+}
  \FunctionTok{coord\_flip}\NormalTok{()}\SpecialCharTok{+}
  \FunctionTok{labs}\NormalTok{(}\AttributeTok{title=}\StringTok{"Observations by sex"}\NormalTok{,}
       \AttributeTok{x=}\StringTok{"sex"}\NormalTok{,}
       \AttributeTok{y=}\ConstantTok{NULL}\NormalTok{) }\CommentTok{\#take away the y axis label}
\end{Highlighting}
\end{Shaded}

\pandocbounded{\includegraphics[keepaspectratio]{midterm_2_files/figure-latex/unnamed-chunk-12-1.pdf}}

\begin{enumerate}
\def\labelenumi{\arabic{enumi}.}
\setcounter{enumi}{3}
\tightlist
\item
  (6 points) How many bears were sampled in each year for both
  populations? Make a line plot(s) to visualize sampling by year for
  both populations. (hint: you will need to pull apart the date to get
  the year in its own column)
\end{enumerate}

\begin{Shaded}
\begin{Highlighting}[]
\FunctionTok{names}\NormalTok{(polar\_bears)}
\end{Highlighting}
\end{Shaded}

\begin{verbatim}
##  [1] "bear_id"       "capture_date"  "popn"          "sex"          
##  [5] "age_class"     "mating"        "denning"       "lat"          
##  [9] "long"          "urea_nitrogen" "creatinine"
\end{verbatim}

\begin{Shaded}
\begin{Highlighting}[]
\NormalTok{polar\_bears}
\end{Highlighting}
\end{Shaded}

\begin{verbatim}
## # A tibble: 519 x 11
##    bear_id capture_date popn  sex   age_class mating denning   lat  long
##      <dbl> <date>       <chr> <chr> <chr>     <chr>  <chr>   <dbl> <dbl>
##  1    6398 1985-04-28   SB    F     A         no     yes      70.4 -148.
##  2    6384 1991-04-13   SB    F     A         no     yes      70.0 -150.
##  3    6520 1992-04-18   SB    F     A         no     yes      70.9 -145.
##  4    6559 2000-04-12   SB    F     A         no     yes      70.0 -142.
##  5   20629 2003-04-16   SB    F     A         no     yes      70.5 -145.
##  6   20876 2007-04-01   SB    F     A         no     yes      71.1 -153.
##  7   20354 2009-04-16   SB    F     A         no     yes      70.7 -145.
##  8   21024 2009-04-25   SB    F     A         no     yes      70.5 -147.
##  9   20764 2009-05-01   SB    F     A         no     yes      70.8 -149.
## 10   21045 2009-05-19   SB    F     A         no     yes      70.7 -148.
## # i 509 more rows
## # i 2 more variables: urea_nitrogen <dbl>, creatinine <dbl>
\end{verbatim}

\begin{Shaded}
\begin{Highlighting}[]
\NormalTok{polar\_bears\_line }\OtherTok{\textless{}{-}}\NormalTok{ polar\_bears }\SpecialCharTok{\%\textgreater{}\%} 
  \FunctionTok{count}\NormalTok{(capture\_date, }\AttributeTok{sort=}\NormalTok{T)}

\NormalTok{polar\_bears\_line}
\end{Highlighting}
\end{Shaded}

\begin{verbatim}
## # A tibble: 377 x 2
##    capture_date     n
##    <date>       <int>
##  1 2008-04-26       4
##  2 2013-04-17       4
##  3 1984-05-06       3
##  4 1989-04-01       3
##  5 1989-04-05       3
##  6 2002-04-19       3
##  7 2004-03-31       3
##  8 2005-04-06       3
##  9 2005-04-12       3
## 10 2006-04-21       3
## # i 367 more rows
\end{verbatim}

4=yrs on its own 6-7=make new variable to store it as the polarbear gets
polarbears bc u need those vairables

\begin{Shaded}
\begin{Highlighting}[]
\NormalTok{polar\_bears\_line }\SpecialCharTok{\%\textgreater{}\%} 
  \FunctionTok{filter}\NormalTok{((capture\_date}\SpecialCharTok{\textgreater{}=}\DecValTok{1980{-}04{-}28}\NormalTok{)) }\SpecialCharTok{\%\textgreater{}\%} 
  \FunctionTok{ggplot}\NormalTok{(}\FunctionTok{aes}\NormalTok{(}\AttributeTok{x=}\NormalTok{capture\_date, }\AttributeTok{y=}\NormalTok{n))}\SpecialCharTok{+}
  \FunctionTok{geom\_line}\NormalTok{()}\SpecialCharTok{+}
  \FunctionTok{geom\_point}\NormalTok{(}\AttributeTok{shape=}\DecValTok{5}\NormalTok{)}\SpecialCharTok{+}
  \FunctionTok{scale\_x\_continuous}\NormalTok{(}\AttributeTok{breaks=}\DecValTok{1985{-}04{-}28}\SpecialCharTok{:}\DecValTok{2016{-}04{-}22}\NormalTok{, }\AttributeTok{minor\_breaks =} \ConstantTok{NULL}\NormalTok{)}\SpecialCharTok{+}
  \FunctionTok{theme}\NormalTok{(}\AttributeTok{axis.text.x =} \FunctionTok{element\_text}\NormalTok{(}\AttributeTok{angle=}\DecValTok{60}\NormalTok{, }\AttributeTok{hjust=}\DecValTok{1}\NormalTok{))}\SpecialCharTok{+}
  \FunctionTok{labs}\NormalTok{(}\AttributeTok{title=}\StringTok{"bears each yr for populations"}\NormalTok{,}
       \AttributeTok{x=}\StringTok{"Year"}\NormalTok{,}
       \AttributeTok{y=}\StringTok{"bears captured"}\NormalTok{)}
\end{Highlighting}
\end{Shaded}

\pandocbounded{\includegraphics[keepaspectratio]{midterm_2_files/figure-latex/unnamed-chunk-16-1.pdf}}

\begin{enumerate}
\def\labelenumi{\arabic{enumi}.}
\setcounter{enumi}{4}
\tightlist
\item
  (2 points) What are the min, median, mean, and max values for
  \texttt{urea\_nitrogen} and \texttt{creatinine}?
\end{enumerate}

\begin{Shaded}
\begin{Highlighting}[]
\FunctionTok{names}\NormalTok{(polar\_bears)}
\end{Highlighting}
\end{Shaded}

\begin{verbatim}
##  [1] "bear_id"       "capture_date"  "popn"          "sex"          
##  [5] "age_class"     "mating"        "denning"       "lat"          
##  [9] "long"          "urea_nitrogen" "creatinine"
\end{verbatim}

\begin{Shaded}
\begin{Highlighting}[]
\NormalTok{polar\_bears }\SpecialCharTok{\%\textgreater{}\%} 
\FunctionTok{summarize}\NormalTok{(}
    \AttributeTok{min\_UR =} \FunctionTok{min}\NormalTok{(urea\_nitrogen, }\AttributeTok{na.rm =} \ConstantTok{TRUE}\NormalTok{),}
    \AttributeTok{median\_UR =} \FunctionTok{median}\NormalTok{(urea\_nitrogen, }\AttributeTok{na.rm =} \ConstantTok{TRUE}\NormalTok{),}
    \AttributeTok{mean\_UR =} \FunctionTok{mean}\NormalTok{(urea\_nitrogen, }\AttributeTok{na.rm =} \ConstantTok{TRUE}\NormalTok{),}
    \AttributeTok{max\_UR =} \FunctionTok{max}\NormalTok{(urea\_nitrogen, }\AttributeTok{na.rm =} \ConstantTok{TRUE}\NormalTok{),}
    \AttributeTok{min\_CR =} \FunctionTok{min}\NormalTok{(creatinine, }\AttributeTok{na.rm =} \ConstantTok{TRUE}\NormalTok{),}
    \AttributeTok{median\_CR =} \FunctionTok{median}\NormalTok{(creatinine, }\AttributeTok{na.rm =} \ConstantTok{TRUE}\NormalTok{),}
    \AttributeTok{mean\_CR =} \FunctionTok{mean}\NormalTok{(creatinine, }\AttributeTok{na.rm =} \ConstantTok{TRUE}\NormalTok{),}
    \AttributeTok{max\_CR =} \FunctionTok{max}\NormalTok{(creatinine, }\AttributeTok{na.rm =} \ConstantTok{TRUE}\NormalTok{),}
    \AttributeTok{.groups =} \StringTok{"drop"}
\NormalTok{  )}
\end{Highlighting}
\end{Shaded}

\begin{verbatim}
## # A tibble: 1 x 8
##   min_UR median_UR mean_UR max_UR min_CR median_CR mean_CR max_CR
##    <dbl>     <dbl>   <dbl>  <dbl>  <dbl>     <dbl>   <dbl>  <dbl>
## 1      1        11    15.1     69   0.35      1.05    1.08   2.35
\end{verbatim}

\begin{enumerate}
\def\labelenumi{\arabic{enumi}.}
\setcounter{enumi}{5}
\tightlist
\item
  (3 points) As explained above, the ratio of \texttt{urea\_nitrogen} to
  \texttt{creatinine} is used as an index of fasting. To calculate the
  ratio, we need to multiply \texttt{urea\_nitrogen} by 0.466, then
  divide by \texttt{creatinine}. Create a new variable called
  \texttt{un\_creat\_ratio} that contains this calculation.
\end{enumerate}

\begin{Shaded}
\begin{Highlighting}[]
\NormalTok{un\_creat\_ratio }\OtherTok{\textless{}{-}}\NormalTok{ polar\_bears }\SpecialCharTok{\%\textgreater{}\%} 
  \FunctionTok{mutate}\NormalTok{(}\AttributeTok{un\_creat\_ratio =}\NormalTok{ (urea\_nitrogen }\SpecialCharTok{*} \FloatTok{0.466}\NormalTok{) }\SpecialCharTok{/}\NormalTok{ creatinine)}

\NormalTok{un\_creat\_ratio}
\end{Highlighting}
\end{Shaded}

\begin{verbatim}
## # A tibble: 519 x 12
##    bear_id capture_date popn  sex   age_class mating denning   lat  long
##      <dbl> <date>       <chr> <chr> <chr>     <chr>  <chr>   <dbl> <dbl>
##  1    6398 1985-04-28   SB    F     A         no     yes      70.4 -148.
##  2    6384 1991-04-13   SB    F     A         no     yes      70.0 -150.
##  3    6520 1992-04-18   SB    F     A         no     yes      70.9 -145.
##  4    6559 2000-04-12   SB    F     A         no     yes      70.0 -142.
##  5   20629 2003-04-16   SB    F     A         no     yes      70.5 -145.
##  6   20876 2007-04-01   SB    F     A         no     yes      71.1 -153.
##  7   20354 2009-04-16   SB    F     A         no     yes      70.7 -145.
##  8   21024 2009-04-25   SB    F     A         no     yes      70.5 -147.
##  9   20764 2009-05-01   SB    F     A         no     yes      70.8 -149.
## 10   21045 2009-05-19   SB    F     A         no     yes      70.7 -148.
## # i 509 more rows
## # i 3 more variables: urea_nitrogen <dbl>, creatinine <dbl>,
## #   un_creat_ratio <dbl>
\end{verbatim}

\begin{enumerate}
\def\labelenumi{\arabic{enumi}.}
\setcounter{enumi}{6}
\tightlist
\item
  (4 points) The authors of the study categorized bears as fasting if
  their \texttt{un\_creat\_ratio} was less than or equal to 10. Use
  \texttt{case\_when()} to create a new variable called \texttt{fasting}
  that categorizes bears as fasting or not.
\end{enumerate}

\begin{Shaded}
\begin{Highlighting}[]
\NormalTok{fasting10 }\OtherTok{\textless{}{-}}\NormalTok{ un\_creat\_ratio }\SpecialCharTok{\%\textgreater{}\%} 
  \FunctionTok{mutate}\NormalTok{(}\AttributeTok{fasting=}\FunctionTok{case\_when}\NormalTok{(un\_creat\_ratio}\SpecialCharTok{\textless{}=}\DecValTok{10} \SpecialCharTok{\textasciitilde{}} \StringTok{"not fasting"}\NormalTok{,}
\NormalTok{                                   un\_creat\_ratio}\SpecialCharTok{\textgreater{}}\DecValTok{10} \SpecialCharTok{\textasciitilde{}} \StringTok{"fasting"}\NormalTok{)) }\SpecialCharTok{\%\textgreater{}\%} 
  \FunctionTok{select}\NormalTok{(popn, un\_creat\_ratio, fasting) }\SpecialCharTok{\%\textgreater{}\%} 
  \FunctionTok{mutate}\NormalTok{(}\AttributeTok{fasting=}\FunctionTok{factor}\NormalTok{(fasting,}
                                \AttributeTok{levels=}\FunctionTok{c}\NormalTok{(}\StringTok{"not fasting"}\NormalTok{,}
                                         \StringTok{"fasting"}\NormalTok{)))}

\NormalTok{fasting10}
\end{Highlighting}
\end{Shaded}

\begin{verbatim}
## # A tibble: 519 x 3
##    popn  un_creat_ratio fasting    
##    <chr>          <dbl> <fct>      
##  1 SB            10.1   fasting    
##  2 SB             0.964 not fasting
##  3 SB            16.4   fasting    
##  4 SB             7.63  not fasting
##  5 SB             4.66  not fasting
##  6 SB             1.79  not fasting
##  7 SB            21.1   fasting    
##  8 SB            19.3   fasting    
##  9 SB             7.57  not fasting
## 10 SB            27.0   fasting    
## # i 509 more rows
\end{verbatim}

\begin{enumerate}
\def\labelenumi{\arabic{enumi}.}
\setcounter{enumi}{7}
\tightlist
\item
  (6 points) The authors found that the number of fasting bears
  increased in the SB population but decreased in the CS population.
  Make a plot(s) that shows the number of fasting bears by year for both
  populations to visualize this trend.
\end{enumerate}

\begin{Shaded}
\begin{Highlighting}[]
\NormalTok{fasting10 }\SpecialCharTok{\%\textgreater{}\%} 
  \FunctionTok{drop\_na}\NormalTok{(fasting) }\SpecialCharTok{\%\textgreater{}\%} \CommentTok{\#get rid of NA\textquotesingle{}s}
  \FunctionTok{ggplot}\NormalTok{(}\FunctionTok{aes}\NormalTok{(}\AttributeTok{x=}\NormalTok{fasting, }\AttributeTok{fill=}\NormalTok{fasting))}\SpecialCharTok{+}
  \FunctionTok{geom\_bar}\NormalTok{()}\SpecialCharTok{+}
  \FunctionTok{labs}\NormalTok{(}\AttributeTok{title=}\StringTok{"fasting bears by year for both populations "}\NormalTok{,}
       \AttributeTok{x=}\StringTok{"fasting bear"}\NormalTok{,}
       \AttributeTok{y=}\ConstantTok{NULL}\NormalTok{)}\SpecialCharTok{+}
  \FunctionTok{theme}\NormalTok{(}\AttributeTok{legend.position=}\StringTok{"none"}\NormalTok{)}
\end{Highlighting}
\end{Shaded}

\pandocbounded{\includegraphics[keepaspectratio]{midterm_2_files/figure-latex/unnamed-chunk-21-1.pdf}}

\begin{enumerate}
\def\labelenumi{\arabic{enumi}.}
\setcounter{enumi}{8}
\tightlist
\item
  (8 points) Where are the polar bears used in this study located? Make
  a static (not leaflet) distribution map that shows the locations of
  bears by population. (hint: use \texttt{color=popn} in the
  \texttt{aes()} of \texttt{geom\_point()})
\end{enumerate}

\begin{Shaded}
\begin{Highlighting}[]
\NormalTok{polar\_bears}
\end{Highlighting}
\end{Shaded}

\begin{verbatim}
## # A tibble: 519 x 11
##    bear_id capture_date popn  sex   age_class mating denning   lat  long
##      <dbl> <date>       <chr> <chr> <chr>     <chr>  <chr>   <dbl> <dbl>
##  1    6398 1985-04-28   SB    F     A         no     yes      70.4 -148.
##  2    6384 1991-04-13   SB    F     A         no     yes      70.0 -150.
##  3    6520 1992-04-18   SB    F     A         no     yes      70.9 -145.
##  4    6559 2000-04-12   SB    F     A         no     yes      70.0 -142.
##  5   20629 2003-04-16   SB    F     A         no     yes      70.5 -145.
##  6   20876 2007-04-01   SB    F     A         no     yes      71.1 -153.
##  7   20354 2009-04-16   SB    F     A         no     yes      70.7 -145.
##  8   21024 2009-04-25   SB    F     A         no     yes      70.5 -147.
##  9   20764 2009-05-01   SB    F     A         no     yes      70.8 -149.
## 10   21045 2009-05-19   SB    F     A         no     yes      70.7 -148.
## # i 509 more rows
## # i 2 more variables: urea_nitrogen <dbl>, creatinine <dbl>
\end{verbatim}

\begin{Shaded}
\begin{Highlighting}[]
\NormalTok{polar\_dups }\OtherTok{\textless{}{-}}\NormalTok{ polar\_bears }\SpecialCharTok{\%\textgreater{}\%} 
  \FunctionTok{distinct}\NormalTok{(popn, }\AttributeTok{.keep\_all =} \ConstantTok{TRUE}\NormalTok{) }
\end{Highlighting}
\end{Shaded}

\begin{Shaded}
\begin{Highlighting}[]
\NormalTok{polar\_dups }\SpecialCharTok{\%\textgreater{}\%} 
  \FunctionTok{select}\NormalTok{(popn) }\SpecialCharTok{\%\textgreater{}\%} 
  \FunctionTok{summary}\NormalTok{()}
\end{Highlighting}
\end{Shaded}

\begin{verbatim}
##      popn          
##  Length:2          
##  Class :character  
##  Mode  :character
\end{verbatim}

\begin{Shaded}
\begin{Highlighting}[]
\CommentTok{\#bbox2 \textless{}{-} make\_bbox(popn, f=0.3)}
\end{Highlighting}
\end{Shaded}

\begin{Shaded}
\begin{Highlighting}[]
\CommentTok{\#map2 \textless{}{-} get\_stadiamap(bbox2, maptype = "stamen\_terrain", zoom=7)}
\CommentTok{\#ggmap(map2)}

\CommentTok{\#ggmap(map2)+}
\CommentTok{\#  geom\_point(data=sharks\_dups, aes(color=popn), size=0.8)+}
\CommentTok{\#  labs(x="Longitude", y="Latitude", title="popn")}
\end{Highlighting}
\end{Shaded}

\begin{enumerate}
\def\labelenumi{\arabic{enumi}.}
\setcounter{enumi}{9}
\tightlist
\item
  (10 points) Build a simple shiny app that shows the range of
  \texttt{un\_creat\_ratio} values for both populations. The user should
  be able to select the fill variable and choose between sex, mating,
  age\_class, denning, and fasting. Please use the code chunk below so
  that your exam knits.
\end{enumerate}

\begin{Shaded}
\begin{Highlighting}[]
\NormalTok{polar\_bears }\OtherTok{\textless{}{-}} 
  \FunctionTok{read\_csv}\NormalTok{(}\StringTok{"data/polarBear\_serumUreaCreatinine\_beaufortChukchi\_rode.csv"}\NormalTok{) }\SpecialCharTok{\%\textgreater{}\%} 
  \FunctionTok{clean\_names}\NormalTok{()}

\NormalTok{ui }\OtherTok{\textless{}{-}} \FunctionTok{dashboardPage}\NormalTok{(}
  
  \FunctionTok{dashboardHeader}\NormalTok{(}\AttributeTok{title =} \StringTok{"range of un\_creat\_ratio values for both populations"}\NormalTok{),}
  
  \FunctionTok{dashboardSidebar}\NormalTok{(}
    
    \FunctionTok{selectInput}\NormalTok{(}\StringTok{"y"}\NormalTok{,}
                \StringTok{"Select Variable"}\NormalTok{,}
                \AttributeTok{choices =} \FunctionTok{c}\NormalTok{(}\StringTok{"sex"}\NormalTok{, }\StringTok{"mating"}\NormalTok{, }\StringTok{"age\_class"}\NormalTok{, }\StringTok{"denning"}\NormalTok{, }\StringTok{"fasting"}\NormalTok{),}
                \AttributeTok{selected =} \StringTok{"popn"}\NormalTok{)}
    
\NormalTok{  ),}
  
  \FunctionTok{dashboardBody}\NormalTok{(}
    
    \FunctionTok{plotOutput}\NormalTok{(}\StringTok{"plot"}\NormalTok{, }\AttributeTok{width =} \StringTok{"600px"}\NormalTok{, }\AttributeTok{height =} \StringTok{"500px"}\NormalTok{)}
    
\NormalTok{  )}
  
\NormalTok{)}

\NormalTok{server }\OtherTok{\textless{}{-}} \ControlFlowTok{function}\NormalTok{(input, output, session) \{}
  
\NormalTok{  output}\SpecialCharTok{$}\NormalTok{plot }\OtherTok{\textless{}{-}} \FunctionTok{renderPlot}\NormalTok{(\{}
    
\NormalTok{    polar\_bears }\SpecialCharTok{\%\textgreater{}\%}
      \FunctionTok{ggplot}\NormalTok{(}\FunctionTok{aes}\NormalTok{(}\AttributeTok{x =}\NormalTok{ un\_creat\_ratio,}
                 \AttributeTok{y =}\NormalTok{ .data[[input}\SpecialCharTok{$}\NormalTok{y]],}
                 \AttributeTok{fill =}\NormalTok{ un\_creat\_ratio)) }\SpecialCharTok{+}
      \FunctionTok{geom\_boxplot}\NormalTok{(}\AttributeTok{alpha =} \FloatTok{0.75}\NormalTok{) }\SpecialCharTok{+}
      \FunctionTok{labs}\NormalTok{(}\AttributeTok{title =} \StringTok{"range of un\_creat\_ratio values by popn CS"}\NormalTok{,}
           \AttributeTok{x =} \StringTok{"CS"}\NormalTok{,}
           \AttributeTok{y =}\NormalTok{ input}\SpecialCharTok{$}\NormalTok{y,}
           \AttributeTok{fill =} \StringTok{"CS"}\NormalTok{) }\SpecialCharTok{+}
      \FunctionTok{theme\_light}\NormalTok{()}
    
\NormalTok{  \})}
  
\NormalTok{\}}

\FunctionTok{shinyApp}\NormalTok{(ui, server)}
\end{Highlighting}
\end{Shaded}

\subsection{Submit the Midterm}\label{submit-the-midterm}

\begin{enumerate}
\def\labelenumi{\arabic{enumi}.}
\tightlist
\item
  Save your work and knit the .rmd file.\\
\item
  Open the .html file and ``print'' it to a .pdf file in Google Chrome
  (not Safari).\\
\item
  Go to the class Canvas page and open Gradescope.\\
\item
  Submit your .pdf file to the midterm assignment- be sure to assign the
  pages to the correct questions.\\
\item
  Commit and push your work to your repository.
\end{enumerate}

\end{document}
